%=============================================================================!
% 9.2  THE VERLET METHOD                                                      !
%=============================================================================!
\section{The Verlet method}
\label{sec:in-verlet}

Verlet's method, which Verlet introduced in 1967 for his simulations of a
liquid of Lennard-Jones atoms\cite{Verlet1967}, gains a whole order over the forward
Euler method at the same cost of one force evaluation per step. This
section derives it from the Taylor series, writes it in the form that
simulations use, and checks its order.

\subsection*{Positions from two Taylor series}

The Taylor series of the motion a step ahead and a step behind are
\begin{align*}
   x(t + \dt) &= x + v\,\dt + \tfrac12 a\,\dt^2 + \tfrac16\dot{a}\,\dt^3
   + \tfrac1{24}\ddot{a}\,\dt^4 + \order{\dt^5} ,\\
   x(t - \dt) &= x - v\,\dt + \tfrac12 a\,\dt^2 - \tfrac16\dot{a}\,\dt^3
   + \tfrac1{24}\ddot{a}\,\dt^4 + \order{\dt^5} ,
\end{align*}
with $x$, $v$, $a$ and the derivatives of $a$ taken at $t$; the second is
the first with $-\dt$ in place of $\dt$, which changes the sign of every
odd power. Adding them, the odd powers cancel, the terms in $\dt^5$
within $\order{\dt^5}$ among them, so the remainder is of order $\dt^6$:
\begin{align*}
   x(t + \dt) + x(t - \dt) &= 2x + a\,\dt^2 + \tfrac1{12}\ddot{a}\,\dt^4
   + \order{\dt^6}
   && \text{$\tfrac1{24} + \tfrac1{24} = \tfrac1{12}$}\\
   x(t + \dt) &= 2x(t) - x(t - \dt) + \dt^2a(t) + \order{\dt^4}
   && \text{subtracting $x(t - \dt)$,}
\end{align*}
the term in $\dt^4$ joining the remainder.
Dropping the last term gives the Störmer--Verlet method,
\begin{equation}
   x_{n+1} = 2x_n - x_{n-1} + \dt^2\,a(x_n) ,
   \label{eq:in-stormer}
\end{equation}
which advances the positions with an error of order $\dt^4$ in one step
from two correct positions, using only the force at the present position.
Its global error is nevertheless of order $\dt^2$, as for velocity Verlet
below. An error made in one position, with the one before it correct,
changes the difference between successive positions, which every later
step carries forward, so over a fixed time each local error grows by a
factor of order $t_{\mathrm f}/\dt$, and the $t_{\mathrm f}/\dt$ of them
give $(t_{\mathrm f}/\dt)^2\times\order{\dt^4} = \order{\dt^2}$. The rule
of \cref{sec:in-steps} that relates the two orders holds only for methods
that step from the present state alone. The velocities do not appear, but
can be recovered by the central difference of \cref{sec:mo-atoms}, $v_n =
(x_{n+1} - x_{n-1})/(2\dt)$, with an error of order $\dt^2$.

\subsection*{Velocity Verlet}

Simulations need the velocities at the same times as the positions, for
the kinetic energy and to start a run from a state. Expanding the
position and the velocity over one step,
\begin{align*}
   x(t + \dt) &= x + v\,\dt + \tfrac12 a\,\dt^2 + \order{\dt^3} ,\\
   v(t + \dt) &= v + a\,\dt + \tfrac12\dot{a}\,\dt^2 + \order{\dt^3} .
\end{align*}
The derivative $\dot{a}$ is not known, but $a(t + \dt) = a + \dot{a}\,\dt
+ \order{\dt^2}$, so $\dot{a}\,\dt = a(t + \dt) - a(t) + \order{\dt^2}$,
and putting this into the velocity series,
\begin{align*}
   v(t + \dt) &= v + a\,\dt + \tfrac12\bigl[a(t + \dt) - a(t)\bigr]\dt +
   \order{\dt^3}\\
   &= v + \tfrac12\bigl[a(t) + a(t + \dt)\bigr]\dt + \order{\dt^3} :
\end{align*}
the velocity changes by the average of the accelerations at the two ends
times the step. Dropping the error terms gives the velocity Verlet
method\cite{Swope1982},
\begin{equation}
   x_{n+1} = x_n + \dt\,v_n + \tfrac12\dt^2a_n , \qquad
   v_{n+1} = v_n + \tfrac12\dt\,(a_n + a_{n+1}) ,
   \label{eq:in-vv}
\end{equation}
with $a_n = a(x_n)$. The method uses $a(x_{n+1})$ in place of the true
$a(t + \dt)$; the two differ by order $\dt^3$, as $x_{n+1}$ and $x(t +
\dt)$ do, which changes $v_{n+1}$ only at order $\dt^4$. Its local error
is of order $\dt^3$ in both the position and the velocity. Written with the velocity at the middle of the
step, $v_{n+\frac12} = v_n + \frac12\dt\,a_n$, it is three moves in turn:
\begin{equation}
   v_{n+\frac12} = v_n + \tfrac12\dt\,a(x_n) , \qquad
   x_{n+1} = x_n + \dt\,v_{n+\frac12} , \qquad
   v_{n+1} = v_{n+\frac12} + \tfrac12\dt\,a(x_{n+1}) ,
   \label{eq:in-kdk}
\end{equation}
a half kick of the velocity, a drift of the position, and another half
kick, in the language of \cref{sec:ha-symplectic}. The force at
$x_{n+1}$ serves both the end of one step and the start of the next, so
each step evaluates the forces once (\cref{fig:in-kdk}).

\begin{figure}[htb]
\centering
\begin{tikzpicture}[line cap=round, line join=round, scale=1.15,
   arr/.style={-{Stealth[length=5pt]}, line width=0.8pt}]
   \draw[black!40] (0,0) -- (8,0);
   \foreach \x/\l in {0/{$t_n$}, 4/{$t_{n+1}$}, 8/{$t_{n+2}$}}
      \draw[black!60] (\x,-0.12) -- (\x,0.12) node[below=7pt, black] {\l};
   \foreach \s in {0,4} {
      \draw[arr, accent] (\s+0.1,0.35) -- (\s+1.9,0.35)
         node[midway, above] {\small kick $\tfrac12\delta t$};
      \draw[arr, black!70] (\s+0.1,0.95) -- (\s+3.9,0.95)
         node[midway, above] {\small drift $\delta t$};
      \draw[arr, accent] (\s+2.1,0.35) -- (\s+3.9,0.35)
         node[midway, above] {\small kick $\tfrac12\delta t$};
   }
   \foreach \x in {0,4,8}
      \node[circle, fill=ochre, inner sep=1.4pt] at (\x,0) {};
   \node[ochre, anchor=west] at (8.25,0.0) {\small forces};
\end{tikzpicture}
\caption{Velocity Verlet as a kick of the velocities over half a step, a
drift of the positions over a whole step with the half-step velocities,
and another half kick. The forces (dots) are evaluated once per step, at
the new positions, and serve the last half kick of one step and the first
of the next.}
\label{fig:in-kdk}
\end{figure}

Started from the same first two positions, velocity Verlet and
Störmer--Verlet produce the same positions after them. From
\cref{eq:in-vv}, two successive steps give
\begin{align*}
   x_{n+1} - x_n &= \dt\,v_n + \tfrac12\dt^2a_n ,\\
   x_n - x_{n-1} &= \dt\,v_{n-1} + \tfrac12\dt^2a_{n-1} ,
\end{align*}
and subtracting the second from the first,
\begin{align*}
   x_{n+1} - 2x_n + x_{n-1} &= \dt\,(v_n - v_{n-1})
   + \tfrac12\dt^2(a_n - a_{n-1})\\
   &= \tfrac12\dt^2(a_{n-1} + a_n) + \tfrac12\dt^2(a_n - a_{n-1})
   && \text{$v_n - v_{n-1} = \tfrac12\dt\,(a_{n-1} + a_n)$}\\
   &= \dt^2a_n && \text{the terms in $a_{n-1}$ cancel,}
\end{align*}
which is \cref{eq:in-stormer}. The leapfrog method, which keeps only the
half-step velocities, $v_{n+\frac12} = v_{n-\frac12} + \dt\,a_n$ and $x_{n+1}
= x_n + \dt\,v_{n+\frac12}$, gives the same positions too (\cref{ex:in-leapfrog}).
Started consistently, the three are one method written three ways, and
the tests of
\texttt{mdlab} confirm that they agree to round-off.

\subsection*{In code}

The function \texttt{velocity\_verlet\_step} of \texttt{mdlab} takes one
step for $N$ atoms:

\begin{code}
def velocity_verlet_step(
    positions: ArrayLike,
    velocities: ArrayLike,
    masses: ArrayLike,
    force: ForceFunction,
    dt: float,
    forces: ArrayLike | None = None,
    force_to_accel: float = FORCE_TO_ACCEL,
) -> tuple[NDArray, NDArray, NDArray]:
    r = np.asarray(positions, dtype=float)
    v = np.asarray(velocities, dtype=float)
    m = np.asarray(masses, dtype=float)
    f = force(r) if forces is None else np.asarray(forces, dtype=float)
    v_half = v + 0.5 * dt * _accel(f, m, force_to_accel)
    r_new = r + dt * v_half
    f_new = force(r_new)
    v_new = v_half + 0.5 * dt * _accel(f_new, m, force_to_accel)
    return r_new, v_new, f_new
\end{code}

\noindent The helper \texttt{\_accel} divides each row of forces by its
atom's mass and multiplies by the factor \texttt{FORCE\_TO\_ACCEL} of
\cref{sec:ne-atoms}. \texttt{ForceFunction} names the kind of argument
\texttt{force} is, a function that takes the positions and returns the
forces. The argument \texttt{forces} may be left out, and then takes the
value \texttt{None}, Python's word for nothing; the line that sets
\texttt{f} then computes the forces at the starting positions, and
otherwise uses those passed in. The forces returned are passed back in
at the next step, so that after the first step \texttt{force} is called
once per step.

\subsection*{Its order, and a claim checked}

In \cref{fig:in-orders}, the global error of velocity Verlet falls along a
line of slope 2.00, so it is a method of order 2, and the error of one
step along a line of slope 2.96: halving the step divides the error at a
given time by 4 and the error of one step by about 8. The introduction of
the TrajCast paper attributes the typical step of 0.5 to \SI{1.0}{\femto
\second} of molecular dynamics to `the cubic scaling of the discretization
error with the integration time step'\cite{Thiemann2026}. That is the
local error of velocity Verlet, of order $\dt^3$, as derived here; over a
fixed time it becomes the global error of order $\dt^2$.

\begin{keyidea}
Adding the Taylor series a step ahead and a step behind cancels the odd
powers and gives the Störmer--Verlet method. Velocity Verlet, the same
method with the velocities kept in step, is a half kick, a drift and a
half kick, costs one force evaluation per step, and is of order 2, with a
local error of order $\dt^3$.
\end{keyidea}

\notebook{09}{measure the order of each method}

\begin{selfcheck}
Show that velocity Verlet follows a body falling under a constant
acceleration $-g$ exactly, whatever the step.
\end{selfcheck}
